\crefalias{section}{appendix}
\section{The Neutron Lifetime Resolution (\texorpdfstring{$\tau_n$}{taun})}

\textbf{The Anomaly:} Precision measurements of the neutron mean lifetime exhibit a $4\sigma$ discrepancy depending on the method used.
\begin{itemize}
    \item \textbf{Bottle Experiments (Storage):} $\tau_{bottle} = 877.75 \pm 0.28$ s.
    \item \textbf{Beam Experiments (Flight):} $\tau_{beam} = 888.0 \pm 2.0$ s.
    \item \textbf{Difference:} $\approx 10$ seconds.
\end{itemize}
Standard physics cannot explain why neutrons appear to live longer when in flight than when confined.

\textbf{The $E_8$ Derivation:} In this framework, the neutron is defined by the \textbf{Symmetry Half-Step} ($M_n \approx M_p + \sigma/2$). Its decay ($n \to p + e + \bar{\nu}$) involves relaxing this symmetry tension through the topological boundary ($\chi$).

We posit that the decay rate is sensitive to the \textbf{Geometric Boundary Condition}:
\begin{enumerate}
    \item \textbf{Confined State (Bottle):} The neutron interacts with the container walls (matter). This reinforces the topological boundary ($\chi$), allowing the decay to proceed at the nominal geometric rate determined by the residual.
    \item \textbf{Free State (Beam):} The neutron propagates in a vacuum. Lacking boundary reinforcement, the decay phase space is dilated by the \textbf{Symmetry-Depth Ratio}. The vacuum impedance suppresses the decay width by the ratio of the Boundary Leakage ($\chi/\sigma$) to the Resonant Depth ($\Delta$).
\end{enumerate}

\textbf{The Geometric Correction Factor:}
We calculate the lifetime dilation factor $\mathcal{D}$ as the resonance depth expanded by the boundary-symmetry ratio:
\begin{equation}
\mathcal{D} = \frac{\Delta}{\Delta - (\chi/\sigma)}
\end{equation}

\textbf{Calculation:}
Using the invariants $\Delta=43$, $\chi=2$, $\sigma=5$:
$$ \chi/\sigma = 0.4 $$
\begin{equation}
\mathcal{D} = \frac{43}{43 - 0.4} = \frac{43}{42.6} \approx \mathbf{1.00939}
\end{equation}

\textbf{Predicted Beam Lifetime:}
Applying this dilation to the precise Bottle measurement:
\begin{equation}
\tau_{beam} = \tau_{bottle} \times \mathcal{D} = 877.75 \text{ s} \times 1.00939
\end{equation}
\begin{equation}
\tau_{beam} \approx \mathbf{886.0 \text{ s}}
\end{equation}

\textbf{Comparison:}
\begin{itemize}
    \item \textbf{Geometric Prediction:} $886.0$ s
    \item \textbf{Experimental Average (Beam):} $888.0 \pm 2.0$ s
    \item \textbf{Result:} The derivation bridges the gap, landing within $1\sigma$ of the beam value.
\end{itemize}

\textbf{Conclusion:} The neutron lifetime discrepancy is not an error; it is a measurement of the \textbf{Topological Environmental Effect}. Bottle experiments measure the intrinsic decay rate; Beam experiments measure the rate dilated by the vacuum's geometric impedance to the symmetry relaxation.
